\chapter{Conclusions}

This comprehensive evaluation of clinical RAG systems provides crucial insights for medical AI system deployment and establishes a foundation for evidence-based model selection in healthcare applications.

\section{Achieved Aims and Objectives}

This study successfully achieved its primary aims of quantifying how different embedding models and small LLMs affect end-to-end retrieval-augmented answering quality, speed, and safety on clinical QA tasks. The comprehensive evaluation of 54 configurations (9 embedding models × 6 LLMs) across 1,080 clinical evaluations provides unprecedented insights into clinical RAG system optimization.

\subsection{Primary Objectives Met}

\begin{enumerate}
    \item \textbf{Comprehensive Model Evaluation}: Successfully evaluated 9 embedding models (biomedbert, e5-base, BioBERT, BioLORD, mini-lm, multi-qa, ms-marco, MedQuAD, mpnet-v2) paired with 6 LLMs (deepseek, gemma, phi3, llama, qwen, tinyllama) across clinical scenarios.

    \item \textbf{Performance Quantification}: Established clear performance metrics with average F1-score of 0.632 ± 0.326, precision of 0.850 ± 0.327, and identified top-performing configurations (multi-qa + phi3: 0.758 F1-score).

    \item \textbf{Clinical Safety Assessment}: Demonstrated high precision across all configurations (0.85+ average) ensuring clinical safety, with specific analysis of hallucination risks and safety indicators.

    \item \textbf{Deployment Guidance}: Developed evidence-based deployment framework for different clinical scenarios (critical care, diagnostic support, general information, high-throughput applications).

    \item \textbf{Efficiency Analysis}: Identified optimal speed-quality trade-offs, demonstrating that clinical RAG systems can achieve both high quality (F1 > 0.72) and reasonable speed (< 66s) without significant compromise.
\end{enumerate}

\subsection{Key Findings and Contributions}

The study revealed several critical insights:

\begin{itemize}
    \item \textbf{Task Specialisation Over Domain Specialisation}: QA-specialised embeddings (multi-qa) outperform medical-domain embeddings, indicating retrieval task alignment is more critical than medical pre-training.
    
    \item \textbf{Model Size Efficiency}: Smaller, well-designed models (1.5-2B parameters) achieve performance competitive with larger models while offering significant efficiency gains.
    
    \item \textbf{Statistical Significance Patterns}: Neither embedding nor LLM selection shows significant main effects (p>0.05), but specific combinations achieve substantial performance differences, emphasizing model combination synergy.
    
    \item \textbf{Clinical Deployment Flexibility}: Multiple configurations achieve clinical-grade performance, providing deployment options based on specific requirements.
\end{itemize}

\section{Critique and Limitations}

\subsection{Current Study Limitations}

\begin{itemize}
    \item \textbf{Dataset Scope}: Limited to MIMIC-IV structure; generalizability to other EHR systems unknown. The evaluation is confined to a specific data format and may not reflect performance across diverse healthcare data systems.
    
    \item \textbf{Evaluation Scale}: 20 questions per configuration; larger question sets would improve statistical power. The current sample size, while comprehensive across models, may not capture all edge cases or rare clinical scenarios.
    
    \item \textbf{Temporal Factors}: Single-time evaluation; performance may vary with model updates and evolving clinical knowledge bases. The rapidly changing landscape of both medical knowledge and AI models requires ongoing evaluation.
    
    \item \textbf{Resource Constraints}: Local hardware limitations may not reflect cloud deployment performance. The evaluation environment may not represent optimal production conditions.
    
    \item \textbf{Methodological Constraints}: The study evaluates on limited MIMIC/synthetic data. Local LLMs (1-4B params) balance privacy/cost but may underperform large models. Results depend on embeddings, chunking, and prompting strategies that may not generalize to all clinical contexts.
\end{itemize}

\subsection{Clinical Deployment Considerations}

While the results demonstrate strong performance across multiple metrics, several factors limit immediate clinical deployment:

\begin{itemize}
    \item \textbf{Regulatory Approval}: The system requires validation through clinical regulatory processes before patient care deployment.
    
    \item \textbf{Integration Complexity}: Real-world EHR integration presents challenges beyond those addressed in this controlled evaluation environment.
    
    \item \textbf{Ethical Considerations}: The balance between AI assistance and human clinical judgment requires careful consideration in practice.
\end{itemize}

\section{Future Work}

The comprehensive analysis reveals specific optimization opportunities and research directions:

\subsection{Technical Enhancements}

\begin{enumerate}
    \item \textbf{Category-Specific Optimization}: Develop specialized retrieval strategies for laboratory data (table-aware chunking), diagnostic queries (ontology-linked indexing), and medication information (structured data extraction).

    \item \textbf{Adaptive Model Selection}: Implement dynamic model switching based on query type and complexity, leveraging the identified performance patterns across clinical categories.

    \item \textbf{Hybrid Architectures}: Combining multiple embedding models for specialized retrieval could address the trade-offs between domain and task specialization.

    \item \textbf{Fine-tuning Studies}: Domain-specific adaptation of pre-trained models could improve performance while maintaining efficiency.
\end{enumerate}

\subsection{Clinical Validation and Deployment}

\begin{enumerate}
    \item \textbf{Healthcare Professional Evaluation}: Clinician-in-the-loop assessment to validate performance in real clinical workflows and decision-making contexts.

    \item \textbf{Real-world Deployment}: Performance monitoring in clinical environments to assess practical utility and identify deployment-specific challenges.

    \item \textbf{Patient Outcome Studies}: Long-term impact assessment on clinical decision-making and patient care outcomes.

    \item \textbf{Multi-institutional Validation}: Evaluate generalizability across different healthcare systems and EHR platforms.
\end{enumerate}

\subsection{System Optimization}

\begin{enumerate}
    \item \textbf{Reliability Enhancement}: Focus on reducing coefficient of variation through ensemble methods or uncertainty quantification, particularly for high-stakes clinical applications.

    \item \textbf{Efficiency Optimization}: Investigate retrieval depth optimization and document filtering strategies to improve the weak correlation between document count and performance quality.

    \item \textbf{Safety and Explainability}: Develop enhanced safety mechanisms and explainable AI features for clinical transparency.
\end{enumerate}
\newpage
\section{Final Remarks}

This research establishes a comprehensive foundation for evidence-based clinical RAG system deployment, providing healthcare organizations with clear guidance for model selection based on specific use cases, performance requirements, and safety considerations. The study demonstrates that clinical RAG systems can achieve high semantic accuracy (0.758 F1-score) and efficient performance (56-66s response time) suitable for healthcare applications.

The key insight that task specialization outweighs domain specialization challenges conventional assumptions about medical AI system design and opens new directions for clinical information retrieval. The strong performance of smaller models provides cost-effective deployment options for resource-constrained healthcare environments while maintaining clinical-grade performance standards.

The comprehensive evaluation framework developed in this work can serve as a template for future clinical AI system evaluations, ensuring systematic, rigorous assessment of medical AI technologies before clinical deployment. The evidence-based deployment recommendations provide immediate practical value for healthcare organizations considering RAG system implementation.

While limitations exist, particularly around generalizability and scale, the research provides substantial evidence that carefully configured RAG systems can meaningfully support clinical information needs while maintaining the safety and accuracy standards required for healthcare applications. The future research directions identified offer clear pathways for continued advancement in clinical AI systems.

This work contributes to the growing body of evidence supporting the responsible deployment of AI in healthcare, demonstrating that with proper evaluation, configuration, and deployment strategies, RAG systems can enhance clinical information access while maintaining the highest standards of patient safety and care quality.
